\chapter{Circuits, advice, and uniformity}
\label{chap:circuits}

This chapter covers Boolean circuits over explicit bases, nonuniform circuit
families and the class $\cc{P/poly}$, its characterization by polynomial
advice, the translation to and from CSLib's circuit model, logspace-uniform
circuit families, the containment $\cc{BPP} \subseteq \cc{P/poly}$, the
small-depth classes $\cc{NC}$, $\cc{AC}$ and $\cc{TC}$, Boolean formulas, and
Barrington's theorem. The library proves that polynomial advice and
polynomial-size circuits give the same class, that logspace-uniform
polynomial-size circuit families decide exactly $\cc{P}$ (Arora--Barak,
Section~6.2), that $\cc{BPP} \subseteq \cc{P/poly}$, the nonuniform Barrington
theorem for variable-bounded logarithmic-depth formulas, and the containment of
$\cc{NC}^1$ in the formula and branching-program classes of that theorem. Still
planned are the interleaving $\cc{NC}^i \subseteq \cc{AC}^i \subseteq
\cc{NC}^{i+1}$, the containment of logarithmic-depth formulas in $\cc{NC}^1$
(and with it $\cc{NC}^1 = \cc{FormulaNC1}$), and a uniform Barrington theorem,
which is the main open mathematical item. The main open infrastructure items
are a circuit encoding designed for log-space generation and a consolidation of
the library's parallel circuit and formula representations
(Section~\ref{sec:circuits-open}).

Two conventions hold throughout. The size of a circuit counts internal and
output gates but not input wires, and negations on gate inputs are free; this
differs from Arora--Barak, who count input vertices and NOT gates, but the
difference does not affect polynomial size bounds. A circuit family stores an
explicit answer for the empty input, because the circuit model requires at
least one input wire. Every class in this chapter is nonuniform unless its name
says otherwise.

\section{The circuit model}
\label{sec:circuits-model}

\begin{definition}[Boolean circuit]
  \label{def:circuit}
  \lean{Complexity.Basis, Complexity.Gate, Complexity.Gate.eval,
    Complexity.Circuit, Complexity.Circuit.eval, Complexity.Circuit.size,
    Complexity.Circuit.depth}
  \leanok
  \uses{def:bitstring}
  A \emph{basis} is a type of operations, each with an arity constraint (any
  fan-in, fan-in exactly $k$, or fan-in at most $k$) and an evaluation map on
  bit strings of admissible length. A circuit over a basis $B$ with $N \ge 1$
  inputs, $M \ge 1$ outputs and $G$ internal gates consists of $G$ internal
  gates and $M$ output gates. A gate applies an operation of $B$ to a list of
  wires whose length the operation admits, and each gate input carries its own
  negation flag. The wires are the $N$ inputs followed by the $G$ internal
  gates. Internal gate $i$ may read only wires $0, \dots, N + i - 1$, so the
  circuit is acyclic by construction, and an output gate may read any wire.
  Evaluation $C : \{0,1\}^N \to \{0,1\}^M$ gives each input wire its input bit
  and each gate the value of its operation on its input wires, each negated
  when its flag is set. The \emph{size} of $C$ is $G + M$. The depth of an
  input wire is $0$, the depth of a gate is one more than the largest depth of
  a wire it reads ($1$ for a gate with no inputs), and the \emph{depth} of $C$
  is the largest depth of an output gate.
\end{definition}

\begin{definition}[AND/OR bases]
  \label{def:circuits-bases}
  \lean{Complexity.Basis.andOr2, Complexity.Basis.boundedAndOr,
    Complexity.Basis.unboundedAndOr}
  \leanok
  \uses{def:circuit}
  The bases whose operations are AND and OR with fan-in exactly $2$, with
  fan-in at most $k$, and with any fan-in. In the last two a gate may have no
  inputs, and then computes the empty conjunction $1$ or the empty disjunction
  $0$. Negation is available only through the free per-input flags. Unless
  stated otherwise, a circuit is a circuit over the fan-in-two basis.
\end{definition}

\begin{definition}[Circuit size complexity]
  \label{def:circuits-size-complexity}
  \lean{Complexity.CompleteBasis, Complexity.instCompleteBasisAndOr2,
    Complexity.Circuit.Realizable, Complexity.Circuit.sizeComplexityWithTop,
    Complexity.Circuit.sizeComplexityWithTop_eq_top_iff,
    Complexity.Circuit.sizeComplexity, Complexity.Circuit.sizeComplexityWithTop_eq_coe}
  \leanok
  \uses{def:circuit, def:circuits-bases}
  A basis is \emph{complete} if, for all $N, M \ge 1$, every function
  $\{0,1\}^N \to \{0,1\}^M$ is computed by some circuit over it; the fan-in-two
  AND/OR basis is complete. For $N \ge 1$, $f : \{0,1\}^N \to \{0,1\}$ and a
  basis $B$, $\mathrm{size}_B(f)$ is the least size of a single-output circuit
  over $B$ computing $f$. The general version takes values in
  $\mathbb{N} \cup \{\infty\}$ and equals $\infty$ exactly when no circuit over
  $B$ computes $f$. For a complete basis a natural-number-valued version is
  provided, and it agrees with the general one.
\end{definition}

\begin{lemma}[Dependency graph]
  \label{lem:circuits-dependency-graph}
  \lean{Complexity.Circuit.dependencyGraph, Complexity.Circuit.totalFanIn,
    Complexity.Circuit.dependencyGraph_adj_val_lt,
    Complexity.Circuit.dependencyGraph_isAcyclic,
    Complexity.Circuit.dependencyGraph_edge_card_le_totalFanIn}
  \leanok
  \uses{def:circuit}
  The wiring digraph of a circuit has its $N + G + M$ input, internal and
  output vertices, numbered in that order, and an edge from each wire to every
  gate that reads it. Every edge goes from a smaller index to a larger one, so
  the digraph is acyclic, and its number of edges is at most the total fan-in
  of the circuit (the number of gate-input occurrences, output gates included).
  The bound can be strict, since a gate that reads one wire twice contributes
  one edge.
\end{lemma}
\begin{proof}
  \leanok
  Internal gate $i$ reads only wires below $N + i$, and the output vertices
  come last. Each edge comes from at least one gate-input occurrence.
\end{proof}

\begin{lemma}[Hardwiring an input prefix]
  \label{lem:circuits-hardwiring}
  \lean{Complexity.Circuit.restrictPrefix, Complexity.Circuit.restrictPrefix_eval,
    Complexity.Circuit.restrictPrefix_size}
  \leanok
  \uses{def:circuit}
  Let $C$ be a fan-in-two circuit on $k + m$ inputs with $m \ge 1$, $M \ge 1$
  outputs and $G$ internal gates, and let $s \in \{0,1\}^k$. There is a
  fan-in-two circuit $C_s$ on $m$ inputs, again with $M$ outputs and $G$
  internal gates and hence of the same size as $C$, such that $C_s(y) = C(s\,y)$
  for every $y \in \{0,1\}^m$, where $s\,y$ is the concatenation.
\end{lemma}
\begin{proof}
  \leanok
  Fix the inputs one at a time; each step removes one input wire, rewires the
  gates that read it so that they compute the same values, and keeps the
  gates.
\end{proof}

\begin{lemma}[Existential quantification of inputs]
  \label{lem:circuits-exists-quantify}
  \lean{Complexity.existsQuantify, Complexity.sizeComplexity_restrictFirst_le,
    Complexity.sizeComplexity_or_le, Complexity.sizeComplexity_existsQuantify_le}
  \leanok
  \uses{def:circuits-size-complexity}
  Write $\mathrm{size}$ for the fan-in-two AND/OR size complexity. Let
  $m \ge 1$, let $f : \{0,1\}^{k + m} \to \{0,1\}$, and let
  $g : \{0,1\}^m \to \{0,1\}$ output $1$ on $y$ exactly when $f(x\,y) = 1$ for
  some $x \in \{0,1\}^k$. Then
  $\mathrm{size}(g) \le 2^k (\mathrm{size}(f) + 1)$.
\end{lemma}
\begin{proof}
  \leanok
  Induction on $k$. Fixing the first input of a function to a constant does
  not increase its size, since the input can be hardwired without adding
  gates, and the disjunction of two functions has size at most the sum of
  their sizes plus one.
\end{proof}

\section{Circuit families and $\cc{P/poly}$}
\label{sec:circuits-families}

\begin{definition}[Circuit family]
  \label{def:circuits-family}
  \lean{Complexity.CircuitFamily, Complexity.CircuitFamily.function,
    Complexity.CircuitFamily.size, Complexity.CircuitFamily.PolynomialSize,
    Complexity.CircuitFamily.language, Complexity.CircuitFamily.Decides}
  \leanok
  \uses{def:circuit, def:language}
  A circuit family over $B$ consists of a single-output circuit $C_n$ over $B$
  on $n$ inputs, with any number of internal gates, for every $n \ge 1$,
  together with an output bit for the empty input. It computes the Boolean
  function family given by $C_n$ at length $n \ge 1$ and by the stored bit at
  length $0$, and it decides the language of strings on which it outputs $1$. Its size at $n \ge 1$ is the size of $C_n$,
  and its size at $0$ is $0$. It has polynomial size if a single polynomial
  $p \in \mathbb{N}[X]$ satisfies $\mathrm{size}(n) \le p(n)$ for every $n$.
\end{definition}

\begin{definition}[$\cc{SIZE}$]
  \label{def:SIZE}
  \lean{Complexity.SIZEWithBasis, Complexity.SIZE}
  \leanok
  \uses{def:circuits-family, def:circuits-bases}
  For $s : \mathbb{N} \to \mathbb{N}$, $\cc{SIZE}(s)$ is the class of languages
  decided by a fan-in-two AND/OR circuit family whose size at every length $n$
  is at most $s(n)$. The same definition is available over any basis.
\end{definition}

\begin{definition}[$\cc{P/poly}$]
  \label{def:PPoly}
  \lean{Complexity.PPoly}
  \leanok
  \uses{def:SIZE}
  $\cc{P/poly} = \bigcup_{p \in \mathbb{N}[X]} \cc{SIZE}(n \mapsto p(n))$.
\end{definition}

\begin{proposition}[Big-O form of $\cc{P/poly}$]
  \label{prop:circuits-ppoly-bigo}
  \lean{Complexity.mem_PPoly_iff}
  \leanok
  \uses{def:PPoly}
  $L \in \cc{P/poly}$ if and only if some fan-in-two circuit family decides $L$
  and has size $O(n^k)$ for some $k$, that is, size at most $C n^k$ at every
  sufficiently large length $n$, for some constant $C$.
\end{proposition}
\begin{proof}
  \leanok
  Convert between polynomial bounds and big-O power bounds.
\end{proof}

\begin{theorem}[Circuits as straight-line programs]
  \label{thm:circuits-straightline}
  \lean{Complexity.Basis.signature, Complexity.Basis.interpretation,
    Complexity.Basis.interpretation_kind,
    Complexity.Circuit.toStraightLine, Complexity.Circuit.eval_toStraightLine,
    Complexity.Circuit.size_toStraightLine, Complexity.Circuit.gatedOutputs_toStraightLine,
    Complexity.Circuit.totalFanIn_toStraightLine}
  \leanok
  \uses{def:circuit, def:circuits-bases}
  Every basis $B$ is a CSLib signature whose operation symbols are the gate
  kinds of $B$: an operation, a fan-in it allows, and a negation flag for each
  input, interpreted exactly as a gate of that kind evaluates. A circuit over
  $B$ with $N \ge 1$ inputs, $G$ internal gates and $M \ge 1$ outputs becomes a
  CSLib circuit over this signature that computes the same function and has
  exactly $G + M$ gates, its size. Every output of the translation is an
  internal gate, and its total fan-in is that of the original circuit. The
  translation goes one way for general bases: the converse, from CSLib
  circuits over $B$'s signature whose outputs are internal gates, is planned.
\end{theorem}
\begin{proof}
  \leanok
  The program lists the internal gates and then the output gates, with each
  wire renumbered in CSLib's layout, and the outputs are the last $M$ gates.
  Induction on the position shows that every gate of the program carries the
  value of the corresponding typed gate.
\end{proof}

\begin{theorem}[CSLib circuits]
  \label{thm:circuits-cslib}
  \lean{Complexity.Circuit.ofCslib, Complexity.Circuit.eval_ofCslib,
    Complexity.Circuit.size_ofCslib, Complexity.Circuit.exists_cslib,
    Complexity.Circuit.sizeComplexity_le_of_cslib,
    Complexity.Circuit.exists_cslib_of_sizeComplexity,
    Complexity.mem_PPoly_iff_cslib, Complexity.PPoly_eq_cslib_PPoly,
    Complexity.SIZE_subset_cslib_SIZE, Complexity.cslib_SIZE_subset_SIZE}
  \leanok
  \uses{def:circuit, def:circuits-bases, def:circuits-size-complexity, def:SIZE,
    def:PPoly}
  Let $N, M \ge 1$. A De Morgan circuit in CSLib's model (constants,
  negation, and binary conjunction and disjunction, every gate counted, and
  outputs designated wires at no cost) with $N$ inputs, $g$ gates, and $M$
  outputs becomes a fan-in-two AND/OR circuit of size exactly $g + M$
  computing the same outputs. Conversely, a fan-in-two AND/OR circuit with $N$
  inputs, $G$ internal gates, and $M$ outputs has a CSLib circuit with at most
  $N + 2G + M$ gates computing the same outputs. For a single function $f$ on
  $N$ bits, a CSLib circuit with $g$ gates computing $f$ gives
  $\mathrm{size}(f) \le g + 1$ for the fan-in-two AND/OR size complexity, and
  $f$ has a CSLib circuit with at most $N + 2\,\mathrm{size}(f)$ gates.
  Consequently $L \in \cc{P/poly}$ if and only if, for some polynomial $p$,
  every length-$n$ slice of $L$, including $n = 0$, is decided by a CSLib
  circuit with at most $p(n)$ gates, and $\cc{P/poly}$ equals CSLib's
  $\cc{P/poly}$. The size classes interleave: $\cc{SIZE}(s)$ lies in
  CSLib's $\cc{SIZE}(n + 2s + 1)$, and CSLib's $\cc{SIZE}(s)$ lies in
  $\cc{SIZE}(s + 1)$.

  CSLib's Boolean $\cc{SIZE}$ and $\cc{P/poly}$ (at most $s(n)$ De Morgan gates
  at every length $n$, and $\bigcup_k \cc{SIZE}(n^k + k)$) are pending CSLib
  work by this library's author, not yet in upstream CSLib; the library builds
  them from the pinned integration branch. The equality of the two
  $\cc{P/poly}$ classes is therefore a consistency check between two
  definitions by the same author, not corroboration by an independently
  reviewed upstream definition.
\end{theorem}
\begin{proof}
  \leanok
  The two models share their wire layout, so the first translation maps each
  CSLib gate to one of ours: $\neg w$ becomes $\neg w \wedge \neg w$ and a
  constant becomes $x_0 \wedge \neg x_0$ or $x_0 \vee \neg x_0$. For the second,
  keep every wire and its negation available, at a cost of $N$ gates for the
  negated inputs, two gates per internal gate, and one gate per output. The
  length-zero slice, which our families answer separately, is a single
  constant gate in CSLib's model.
\end{proof}

\begin{corollary}[A language outside $\cc{P/poly}$]
  \label{cor:circuits-not-ppoly}
  \lean{Complexity.exists_not_mem_PPoly}
  \leanok
  \uses{def:PPoly}
  Some language is not in $\cc{P/poly}$.
\end{corollary}
\begin{proof}
  \leanok
  \uses{thm:circuits-cslib}
  By the equality with CSLib's $\cc{P/poly}$, this is CSLib's theorem that
  some language lies outside its $\cc{P/poly}$: by Shannon's counting argument
  some language has, at every large length $n$, no De Morgan circuit with at
  most $2^n / n$ gates, and every bound $n^k + k$ falls below $2^n / n$ for
  large $n$. That theorem, like CSLib's $\cc{P/poly}$, comes from the pending
  CSLib work described in Theorem~\ref{thm:circuits-cslib}.
\end{proof}

\begin{theorem}[Depth in CSLib's circuit model]
  \label{thm:circuits-cslib-depth}
  \lean{Complexity.Circuit.depth_ofCslib, Complexity.Circuit.exists_cslib_depth_le,
    Complexity.exists_cslib_of_mem_DEPTH, Complexity.mem_DEPTH_of_cslib}
  \leanok
  \uses{thm:circuits-cslib, def:NC-AC}
  Let $N, M \ge 1$, as in Theorem~\ref{thm:circuits-cslib}. In CSLib's model
  inputs have depth $0$, every gate (negations and constants included) adds
  one level, and the depth of a circuit is the largest depth of an output
  wire. The translation of a CSLib De Morgan circuit has depth exactly one
  more, for its output gates. Conversely, a fan-in-two AND/OR circuit with $N$
  inputs, $G$ internal gates, and $M$ outputs has a CSLib circuit computing
  the same outputs with exactly $N + 2G + M$ gates and depth at most one more.
  For the depth classes $\cc{DEPTH}(d)$ of Definition~\ref{def:NC-AC}: if
  $f \in \cc{DEPTH}(d)$, then for every $n \ge 1$ some CSLib circuit of depth
  at most $d(n) + 1$ computes $f_n$; and if for every $n \ge 1$ some CSLib
  circuit of depth at most $d(n)$ computes $f_n$, then
  $f \in \cc{DEPTH}(n \mapsto d(n) + 1)$.
\end{theorem}
\begin{proof}
  \leanok
  \uses{thm:circuits-cslib}
  The first translation replaces each CSLib gate by one gate of ours. For the
  second, keep both rails of every wire: negate the inputs once, compute each
  gate's positive rail directly and its negative rail by De Morgan from the
  complementary rails, so each gate costs two gates and no extra depth.
\end{proof}

\begin{theorem}
  \label{thm:circuits-p-subset-ppoly}
  \lean{Complexity.P_subset_PPoly}
  \leanok
  \uses{def:P, def:PPoly}
  $\cc{P} \subseteq \cc{P/poly}$.
\end{theorem}
\begin{proof}
  \leanok
  \uses{thm:circuits-unrolling-family, prop:circuits-ppoly-bigo}
  Unroll a polynomial-time decider into its circuit family.
\end{proof}

\section{Advice}
\label{sec:circuits-advice}

\begin{definition}[Polynomial advice]
  \label{def:PAdvice}
  \lean{Complexity.Advice, Complexity.PolynomialAdvice, Complexity.advisedInput,
    Complexity.TM.DecidesWithAdviceInTime, Complexity.PAdvice}
  \leanok
  \uses{def:tm, def:pair, def:language}
  An advice function $a : \mathbb{N} \to \{0,1\}^*$ has polynomial length if
  $|a(n)| \le p(n)$ for every $n$ and a polynomial $p \in \mathbb{N}[X]$. A
  deterministic machine $M$ decides $L$ with advice $a$ in time $T$ if on input
  $\langle a(|x|), x \rangle$ (the library's self-delimiting pairing) it halts
  within $T(|x|)$ steps with $1$ in the first output cell when $x \in L$ and
  $0$ otherwise. The time bound is measured in $|x|$, not in the length of the
  paired input. $\cc{PAdvice}$ is the class of languages decided, by a machine
  with any number of work tapes, with polynomial-length advice in time
  $O(n^d)$ for some $d$.
\end{definition}

\begin{theorem}
  \label{thm:circuits-padvice-subset-ppoly}
  \lean{Complexity.PAdvice_subset_PPoly, Complexity.TM.adviceCircuitFamily,
    Complexity.TM.DecidesWithAdviceInTime.adviceCircuitFamily_decides,
    Complexity.TM.adviceCircuitFamily_size_bigO,
    Complexity.TM.DecidesWithAdviceInTime.mem_PPoly}
  \leanok
  \uses{def:PAdvice, def:PPoly}
  $\cc{PAdvice} \subseteq \cc{P/poly}$. More precisely, if a deterministic
  machine decides $L$ with advice $a$ in time $T = O(n^d)$, then hardwiring
  the advice into its acceptance circuits gives a fan-in-two family of size
  $O(n^{3d})$ deciding $L$, whatever the length of the advice.
\end{theorem}
\begin{proof}
  \leanok
  \uses{thm:circuits-acceptance-circuit, lem:circuits-hardwiring}
  At each length, unroll the advised machine and hardwire the fixed advice
  prefix of its input.
\end{proof}

\begin{theorem}
  \label{thm:circuits-ppoly-subset-padvice}
  \lean{Complexity.PPoly_subset_PAdvice,
    Complexity.CircuitFamily.polynomialAdvice_encodeAt,
    Complexity.CircuitFamily.Decides.evalFamilyTM_decidesWithAdviceInTime,
    Complexity.CircuitFamily.adviceEvalTime_bigO,
    Complexity.CircuitFamily.Decides.mem_PAdvice}
  \leanok
  \uses{def:PAdvice, def:PPoly, def:circuits-code}
  $\cc{P/poly} \subseteq \cc{PAdvice}$. More precisely, if a fan-in-two family
  $F$ of size $O(n^d)$ decides $L$, then its family codes
  (Definition~\ref{def:circuits-code}) form polynomial-length advice, with
  which the serialized circuit evaluator of
  Theorem~\ref{thm:circuits-evaluator} decides $L$ in time $O(n^{4(d+1)})$,
  measured in $|x|$.
\end{theorem}
\begin{proof}
  \leanok
  \uses{lem:circuits-codec, thm:circuits-evaluator}
  The advice at length $n$ is the family code at length $n$, which has
  polynomial length by Lemma~\ref{lem:circuits-codec}, and the machine is the
  serialized circuit evaluator, whose quadratic running time in the length of
  the paired input is polynomial in $|x|$.
\end{proof}

\begin{corollary}
  \label{cor:circuits-padvice-eq-ppoly}
  \lean{Complexity.PAdvice_eq_PPoly}
  \leanok
  \uses{def:PAdvice, def:PPoly}
  $\cc{PAdvice} = \cc{P/poly}$.
\end{corollary}
\begin{proof}
  \leanok
  \uses{thm:circuits-padvice-subset-ppoly, thm:circuits-ppoly-subset-padvice}
  Combine the two containments.
\end{proof}

\section{Serialized circuits and their evaluation}
\label{sec:circuits-codes}

\begin{definition}[Circuit and family codes]
  \label{def:circuits-code}
  \lean{Complexity.CircuitCode.RawCircuit, Complexity.CircuitCode.encodeCircuit,
    Complexity.CircuitCode.evalCode, Complexity.CircuitFamily.encodeAt,
    Complexity.CircuitCode.evalFamilyCode,
    Complexity.CircuitCode.circuitEvalLanguage}
  \leanok
  \uses{def:circuits-family, def:pair}
  A fan-in-two single-output circuit is serialized as its gate count in
  terminated unary followed by its gates in index order, the output gate last.
  A gate is written as an operation bit, two negation flags, and two absolute
  wire references in terminated unary. The partial evaluator
  $\mathrm{evalCode}(N, c, x)$ decodes $c$, rejects inputs of the wrong length
  and codes that are malformed, have trailing bits, are empty, or reference a
  wire that is not yet computed, and evaluates the gates in order with a memo
  array. The code of a family at length $0$ is the bit
  $0$ followed by the stored answer; at length $n \ge 1$ it is the bit $1$
  followed by the code of $C_n$. The family evaluator runs a code $c$ on a
  string $x$ as follows: if $x$ is empty, $c$ must be $0\,b$ for a bit $b$,
  and the answer is $b$; if $x$ is nonempty, $c$ must be $1\,c'$, and the
  answer is $\mathrm{evalCode}(|x|, c', x)$; every other case is rejected. The
  evaluation language is the set of strings $\langle c, x \rangle$ (the
  library's self-delimiting pairing) on which the family evaluator returns
  $1$.
\end{definition}

\begin{lemma}[Codec correctness]
  \label{lem:circuits-codec}
  \lean{Complexity.CircuitCode.evalCode_encodeCircuit,
    Complexity.CircuitCode.encodeCircuit_length_le_size}
  \leanok
  \uses{def:circuits-code}
  For a fan-in-two single-output circuit $C$ with $N \ge 1$ inputs and size $s$,
  $\mathrm{evalCode}(N, \mathrm{code}(C), x) = C(x)$ for every
  $x \in \{0,1\}^N$, and $|\mathrm{code}(C)| \le 1 + s\,(2(N + s) + 6)$.
\end{lemma}
\begin{proof}
  \leanok
  The iterative evaluator agrees gate by gate with typed evaluation.
\end{proof}

\begin{theorem}[Serialized evaluator]
  \label{thm:circuits-evaluator}
  \lean{Complexity.CircuitCode.Machine.evalFamilyTM,
    Complexity.CircuitCode.Machine.evalFamilyTime,
    Complexity.CircuitCode.Machine.evalFamilyTM_decidesInTime,
    Complexity.CircuitCode.Machine.evalFamilyTime_bigO_quadratic,
    Complexity.circuitEvalLanguage_mem_P}
  \leanok
  \uses{def:circuits-code, def:tm-decides, def:P}
  A fixed deterministic machine with three work tapes decides the evaluation
  language within $4n + 20(n + 1)^2 + 17 = O(n^2)$ steps on every input of
  length $n$; malformed pairs and malformed codes are rejected. In particular
  the evaluation language is in $\cc{P}$.
\end{theorem}
\begin{proof}
  \leanok
  Validate the outer pair and copy the code and the input onto work tapes in
  linear time. Then stream the gates once, appending each gate value to a memo
  tape and reading the two referenced values by unary lookup; each gate costs
  $O(|\mathrm{memo}|)$ steps, for a total of $20(|c| + |x| + 1)^2$ steps in
  the core.
\end{proof}

\section{Functional unrolling}
\label{sec:circuits-unrolling}

The circuits below compute the output bit of a bounded run. They are not the
Cook--Levin CNF, which expresses the existence of an accepting tableau rather
than computing the output of a deterministic run.

\begin{theorem}[Acceptance circuit of a bounded run]
  \label{thm:circuits-acceptance-circuit}
  \lean{Complexity.CircuitUnrolling.acceptanceCircuit,
    Complexity.CircuitUnrolling.acceptanceCircuit_eval,
    Complexity.CircuitUnrolling.acceptanceSizeCoeff,
    Complexity.CircuitUnrolling.acceptanceCircuit_size_le,
    Complexity.CircuitUnrolling.canonicalAcceptanceCircuit,
    Complexity.CircuitUnrolling.card_acceptingChoices_eq_acceptCount}
  \leanok
  \uses{def:ntm, def:circuit}
  Let $M$ be a nondeterministic machine and $T, n, m \in \mathbb{N}$ with
  $m \ge 1$. For every assignment of the $n$ data positions and the $T$ choice
  positions to input wires among $m$ inputs, there is a fan-in-two circuit on
  $m$ inputs that, on every input whose data wires carry $x \in \{0,1\}^n$ and
  whose choice wires carry $c \in \{0,1\}^T$, outputs $1$ exactly when the
  $T$-step trace of $M$ on $x$ under $c$ is in the halting state with $1$ in
  the first output cell. Its size is at most $\kappa_M (T + 2)^3$ for a
  constant $\kappa_M$ depending only on $M$, not on $n$ or $m$. With the
  choices placed first ($m = T + n \ge 1$), the number of $c \in \{0,1\}^T$
  on which the circuit accepts $c\,x$ equals the number of accepting choice
  sequences of length $T$ of $M$ on $x$.
\end{theorem}
\begin{proof}
  \leanok
  Encode bounded configurations one-hot, compile the initial configuration and
  one transition layer as Boolean formula fragments, tile $T$ layers, and
  append a final halt-and-output gate.
\end{proof}

\begin{theorem}[Unrolling family]
  \label{thm:circuits-unrolling-family}
  \lean{Complexity.TM.unrollingCircuitFamily,
    Complexity.TM.DecidesInTime.unrollingCircuitFamily_decides,
    Complexity.TM.unrollingCircuitFamily_size_bigO}
  \leanok
  \uses{def:tm-decides, def:circuits-family}
  For a deterministic machine $M$ and $f : \mathbb{N} \to \mathbb{N}$, the
  unrolling family has at length $n \ge 1$ the acceptance circuit of $M$ at
  horizon $f(n)$ with every choice input hardwired to $0$, and at length $0$
  the output bit of the $f(0)$-step run on the empty input. If $M$ decides $L$
  in time $f$, the family decides $L$. If $f = O(n^d)$, the family has size
  $O(n^{3d})$.
\end{theorem}
\begin{proof}
  \leanok
  \uses{thm:circuits-acceptance-circuit, lem:circuits-hardwiring}
  A deterministic machine ignores its choices, so any fixed choice string
  gives the output bit.
\end{proof}

\section{Logspace uniformity}
\label{sec:circuits-uniform}

\begin{definition}[Logspace-uniform $\cc{P/poly}$]
  \label{def:UniformPPoly}
  \lean{Complexity.unaryList, Complexity.CircuitFamily.Uniform,
    Complexity.UniformPPoly}
  \leanok
  \uses{def:PPoly, def:circuits-code, def:FL}
  A fan-in-two circuit family is logspace-uniform if some function in
  $\cc{FL}$ maps $1^n$ to the family code at length $n$, for every $n$, as in
  Arora--Barak, Section~6.2. The generator reads $n$ in unary, and its
  output at $n = 0$ carries the stored empty-input answer. The class
  $\cc{UniformPPoly}$ consists of the languages decided by a logspace-uniform
  circuit family of polynomial size. P-uniformity is not used.
\end{definition}

\begin{proposition}
  \label{prop:circuits-uniform-subset-ppoly}
  \lean{Complexity.UniformPPoly_subset_PPoly}
  \leanok
  \uses{def:UniformPPoly}
  $\cc{UniformPPoly} \subseteq \cc{P/poly}$.
\end{proposition}
\begin{proof}
  \leanok
  Forget the generator.
\end{proof}

\begin{theorem}
  \label{thm:circuits-uniform-subset-p}
  \lean{Complexity.UniformPPoly_subset_P}
  \leanok
  \uses{def:UniformPPoly, def:P}
  $\cc{UniformPPoly} \subseteq \cc{P}$.
\end{theorem}
\begin{proof}
  \leanok
  \uses{thm:circuits-evaluator, thm:space-L-subset-P, thm:polytime-preimage}
  Since $\cc{FL} \subseteq \cc{FP}$, the map
  $x \mapsto \langle \mathrm{gen}(1^{|x|}), x \rangle$ is in $\cc{FP}$. The
  language is the preimage of the evaluation language under this map, and
  $\cc{P}$ is closed under preimages by $\cc{FP}$ functions. The size bound of
  the family is not used.
\end{proof}

\begin{lemma}[Padded tableau circuits and their code map]
  \label{lem:circuits-padded-code-fl}
  \lean{Complexity.TM.directUnrollingCircuitFamily,
    Complexity.TM.directUnrollingGateBound,
    Complexity.TM.paddedDirectUnrollingRawCircuit,
    Complexity.TM.paddedDirectUnrollingCircuitFamily,
    Complexity.TM.paddedDirectUnrollingCode,
    Complexity.TM.paddedDirectUnrollingCircuitFamily_encodeAt,
    Complexity.TM.DecidesInTime.paddedDirectUnrollingCircuitFamily_decides,
    Complexity.TM.paddedDirectUnrollingCircuitFamily_size_bigO,
    Complexity.TM.directSerializerHorizonPolynomial,
    Complexity.TM.paddedDirectUnrollingCode_mem_FL}
  \leanok
  \uses{def:tm, def:tm-decides, def:circuits-code, def:FL,
    thm:circuits-acceptance-circuit}
  Let $M$ be a deterministic machine and $h : \mathbb{N} \to \mathbb{N}$. The
  direct unrolling of $M$ at horizon $h$ has at length $n \ge 1$ the acceptance
  circuit of Theorem~\ref{thm:circuits-acceptance-circuit} at horizon $h(n)$ on
  $n$ inputs, the data on the inputs in order and every choice position on
  input $0$. The padded direct unrolling $P_{M,h}$ appends to this gate list
  constant-$0$ gates up to $\kappa_M (h(n) + 2)^3$ gates and then a gate
  copying the original output, so its gate count is a closed-form cubic in
  $h(n)$; at length $0$ it stores the output bit of the $h(0)$-step run on the
  empty input. Let $D_{M,h}(n)$ be the family code of $P_{M,h}$ at length $n$.
  If $M$ decides $L$ in time $h$, then $P_{M,h}$ decides $L$, and if
  $h = O(n^d)$, then $P_{M,h}$ has size $O(n^{3d})$. For every polynomial
  $q \in \mathbb{N}[X]$ and $h(n) = q(n) + n + 1$, the map
  $x \mapsto D_{M,h}(|x|)$ is in $\cc{FL}$.
\end{lemma}
\begin{proof}
  \leanok
  The padding gates are never read and the final gate copies the acceptance
  bit, so correctness and the size bound come from
  Theorem~\ref{thm:circuits-acceptance-circuit}, the choices being irrelevant
  to a deterministic machine. For the code map, the direct unrolling is an
  initialization fragment, a stream of fixed-size
  transition fragments whose wire offsets have closed forms, and a final gate.
  A streaming transducer computes the offsets with binary counters, addition,
  multiply-add and fixed-polynomial evaluation, and writes each reference in
  terminated unary. One polynomial bounds every counter value, so the
  auxiliary space is logarithmic.
\end{proof}

\begin{theorem}
  \label{thm:circuits-p-subset-uniform}
  \lean{Complexity.P_subset_UniformPPoly, Complexity.TM.DecidesInTime.mem_UniformPPoly,
    Complexity.TM.DecidesInTime.directSerializerHorizon}
  \leanok
  \uses{def:P, def:UniformPPoly}
  $\cc{P} \subseteq \cc{UniformPPoly}$. More precisely, every language decided
  by a deterministic machine within $q(n)$ steps, for a polynomial
  $q \in \mathbb{N}[X]$, is in $\cc{UniformPPoly}$.
\end{theorem}
\begin{proof}
  \leanok
  \uses{lem:circuits-padded-code-fl}
  A machine that decides $L$ in time $q$ also decides it in time
  $h(n) = q(n) + n + 1$. At horizon $h$ its padded direct unrolling decides
  $L$, has polynomial size, and has a code map in $\cc{FL}$.
\end{proof}

\begin{theorem}[Logspace-uniform circuits capture $\cc{P}$]
  \label{thm:circuits-uniform-eq-p}
  \lean{Complexity.UniformPPoly_eq_P}
  \leanok
  \uses{def:P, def:UniformPPoly}
  $\cc{UniformPPoly} = \cc{P}$.
\end{theorem}
\begin{proof}
  \leanok
  \uses{thm:circuits-uniform-subset-p, thm:circuits-p-subset-uniform}
  Combine the two containments.
\end{proof}

\section{$\cc{BPP}$ is contained in $\cc{P/poly}$}
\label{sec:circuits-bpp}

\begin{lemma}[One seed for all inputs of a length]
  \label{lem:circuits-uniform-seed}
  \lean{Complexity.NTM.uniformSeedRuns, Complexity.NTM.exists_uniform_correct_seed}
  \leanok
  \uses{def:ntm}
  Let $M$ be a probabilistic machine and $f : \mathbb{N} \to \mathbb{N}$ such
  that, for every $x$, the fraction of choice strings $c \in \{0,1\}^{f(|x|)}$
  on which the $f(|x|)$-step trace of $M$ on $x$ under $c$ is halted with
  output $1$ is at least $2/3$ if $x \in L$ and at most $1/3$ if
  $x \notin L$; no halting assumption is made. For every $n$ there is a
  single string of $r(n) f(n)$ random bits, where $r(n) = 12(n + 1) + 1$, on
  which the strict majority vote of these $f(n)$-step runs of $M$, one on each
  of the $r(n)$ consecutive blocks of $f(n)$ bits, is correct for every
  $x \in \{0,1\}^n$.
\end{lemma}
\begin{proof}
  \leanok
  Amplification makes the error on each input less than $2^{-(n+1)}$, and a
  union bound over the $2^n$ inputs leaves a good seed.
\end{proof}

\begin{theorem}[Hardwired amplification]
  \label{thm:circuits-hardwired-amplification}
  \lean{Complexity.NTM.uniformCorrectSeed, Complexity.NTM.hardwiredAmplificationFamily,
    Complexity.NTM.hardwiredAmplificationFamily_decides,
    Complexity.NTM.hardwiredAmplificationFamily_size_bigO}
  \leanok
  \uses{def:ntm, def:circuits-family, lem:circuits-uniform-seed}
  Under the hypotheses of Lemma~\ref{lem:circuits-uniform-seed}, the family
  whose circuit at length $n \ge 1$ hardwires a seed selected by that lemma
  into $r(n)$ parallel copies of the acceptance circuit at horizon $f(n)$,
  followed by a strict-majority fragment, and whose answer at length $0$ is
  whether the empty string is in $L$, is a fan-in-two circuit family deciding
  $L$. If $f = O(n^d)$, its size is $O(n^{3d+4})$.
\end{theorem}
\begin{proof}
  \leanok
  \uses{lem:circuits-uniform-seed, thm:circuits-acceptance-circuit,
    lem:circuits-hardwiring}
  Compose copies of the original acceptance circuit rather than unrolling the
  repeated machine, whose control state would store the whole vote vector.
\end{proof}

\begin{theorem}[Adleman]
  \label{thm:circuits-bpp-subset-ppoly}
  \lean{Complexity.BPP_subset_PPoly}
  \leanok
  \uses{def:BPP, def:PPoly}
  $\cc{BPP} \subseteq \cc{P/poly}$.
\end{theorem}
\begin{proof}
  \leanok
  \uses{thm:circuits-hardwired-amplification, prop:circuits-ppoly-bigo}
  Apply the hardwired amplification family to a $\cc{BPP}$ machine.
\end{proof}

\begin{corollary}
  \label{cor:circuits-bpp-subset-padvice}
  \lean{Complexity.BPP_subset_PAdvice}
  \leanok
  \uses{def:BPP, def:PAdvice}
  $\cc{BPP} \subseteq \cc{PAdvice}$.
\end{corollary}
\begin{proof}
  \leanok
  \uses{thm:circuits-bpp-subset-ppoly, thm:circuits-ppoly-subset-padvice}
  Compose the two containments.
\end{proof}

\section{Small-depth classes}
\label{sec:circuits-depth}

\begin{definition}[$\cc{DEPTH}$, $\cc{NC}^i$ and $\cc{AC}^i$]
  \label{def:NC-AC}
  \lean{Complexity.DEPTHWithBasis, Complexity.DEPTH, Complexity.polylogDepth,
    Complexity.NC, Complexity.AC, Complexity.NC0, Complexity.NC1, Complexity.AC0}
  \leanok
  \uses{def:circuits-family, def:circuits-bases}
  These are classes of Boolean function families, one function
  $\{0,1\}^n \to \{0,1\}$ for each $n$ (including $n = 0$), not of languages.
  The depth of a circuit family at length $n \ge 1$ is the depth of $C_n$, and
  at length $0$ it is $0$. For $d : \mathbb{N} \to \mathbb{N}$,
  $\cc{DEPTH}(d)$ consists of the families computed by a fan-in-two AND/OR
  circuit family of depth at most $d(n)$ at every length $n$, with no size
  bound; the same definition is available over any basis. A family is in
  $\cc{NC}^i$ if a fan-in-two AND/OR circuit family of polynomial size and
  depth at most $c\,(\lfloor \log_2 n \rfloor + 1)^i$ at every length $n$, for
  some constant $c \in \mathbb{N}$ (with $\lfloor \log_2 0 \rfloor = 0$),
  computes it. $\cc{AC}^i$ is defined the same way with unbounded fan-in.
  $\cc{NC}^0$, $\cc{NC}^1$ and $\cc{AC}^0$ are the special cases. These classes
  are nonuniform.
\end{definition}

\begin{definition}[$\cc{TC}^i$]
  \label{def:circuits-TC}
  \lean{Complexity.ThresholdOp, Complexity.Basis.threshold, Complexity.TC,
    Complexity.TC0}
  \leanok
  \uses{def:NC-AC}
  A threshold gate with cutoff $t$ outputs $1$ when at least $t$ of its
  inputs are $1$; the threshold basis has unweighted threshold gates of any
  fan-in, with the free negation flags on their inputs. $\cc{TC}^i$ is defined
  like $\cc{AC}^i$ over this basis.
\end{definition}

\begin{proposition}[Depth hierarchies]
  \label{prop:circuits-depth-classes}
  \lean{Complexity.NC_mono, Complexity.AC_mono, Complexity.TC_mono,
    Complexity.NC0_subset_NC1, Complexity.AC_subset_TC, Complexity.AC0_subset_TC0}
  \leanok
  \uses{def:NC-AC, def:circuits-TC}
  For $i \le j$, $\cc{NC}^i \subseteq \cc{NC}^j$, $\cc{AC}^i \subseteq \cc{AC}^j$
  and $\cc{TC}^i \subseteq \cc{TC}^j$; in particular
  $\cc{NC}^0 \subseteq \cc{NC}^1$. For every $i$,
  $\cc{AC}^i \subseteq \cc{TC}^i$; in particular
  $\cc{AC}^0 \subseteq \cc{TC}^0$.
\end{proposition}
\begin{proof}
  \leanok
  Monotonicity is monotonicity of the depth envelope. For
  $\cc{AC}^i \subseteq \cc{TC}^i$, transport circuits along the embedding of
  AND and OR gates as threshold gates, which preserves semantics, size and
  depth exactly.
\end{proof}

\begin{proposition}[Interleaving]
  \label{prop:circuits-nc-ac-interleave}
  \uses{def:NC-AC}
  For every $i$, $\cc{NC}^i \subseteq \cc{AC}^i \subseteq \cc{NC}^{i+1}$.
\end{proposition}
\begin{proof}
  A fan-in-two gate is an unbounded gate. Conversely, size does not count
  wires, so first merge repeated inputs of each unbounded gate: a gate of a
  circuit with $n$ inputs and $G$ internal gates then reads at most
  $2(n + G)$ distinct literals, a polynomial number. Replace each such gate by
  a balanced fan-in-two tree of depth $O(\log n)$, and a gate with no inputs
  by a constant built from an input and its negation.
\end{proof}

\section{Formulas and Barrington's theorem}
\label{sec:circuits-barrington}

\begin{definition}[Boolean formula]
  \label{def:circuits-formula}
  \lean{Complexity.BoolFormula, Complexity.BoolFormula.eval,
    Complexity.BoolFormula.size, Complexity.BoolFormula.depth,
    Complexity.BoolFormula.vars}
  \leanok
  A Boolean formula is a tree built from variables $x_i$ ($i \in \mathbb{N}$),
  the two constants, negation, and binary conjunction and disjunction. It is
  evaluated under a total assignment $\mathbb{N} \to \{0,1\}$. Its size is its
  number of nodes, its depth is the length of a longest root-to-leaf path,
  with each negation counted as one level, and its variables are the indices
  $i$ of the variables $x_i$ occurring in it.
\end{definition}

\begin{lemma}[Unfolding a circuit output]
  \label{lem:circuits-output-formula}
  \lean{Complexity.Circuit.outputFormula, Complexity.Circuit.eval_outputFormula,
    Complexity.Circuit.vars_outputFormula_lt,
    Complexity.Circuit.depth_outputFormula_le_outputDepth}
  \leanok
  \uses{def:circuit, def:circuits-formula}
  Unfolding one output of a fan-in-two circuit $C$ with $N \ge 1$ inputs gives
  a formula whose variables are all below $N$, which under every assignment
  $\alpha$ computes that output of $C$ on
  $(\alpha(0), \dots, \alpha(N - 1))$, and whose depth is at most twice the
  depth of that output gate. No size bound is claimed, since shared
  subcircuits are duplicated.
\end{lemma}
\begin{proof}
  \leanok
  Structural recursion on wire indices; an edge negation followed by its gate
  costs two formula levels.
\end{proof}

\begin{definition}[Permutation branching program]
  \label{def:circuits-bp}
  \lean{Complexity.BPInstr, Complexity.BP, Complexity.BP.eval,
    Complexity.BP.vars, Complexity.BP.Computes}
  \leanok
  A width-$w$ permutation branching program is a list of instructions
  $(i, \pi_0, \pi_1)$ with $i \in \mathbb{N}$ and $\pi_0, \pi_1 \in S_w$; the
  instruction reads variable $i$. Under an assignment $\alpha$, an instruction
  selects $\pi_{\alpha(i)}$, and the program evaluates to the product
  $\pi^{(1)} \pi^{(2)} \cdots \pi^{(k)}$ of the selected permutations in
  $S_w$, where $(\pi \tau)(y) = \pi(\tau(y))$. A program computes $f$ through
  $\sigma$ if, under every assignment $\alpha$, it evaluates to $\sigma$ when
  $f(\alpha) = 1$ and to the identity otherwise.
\end{definition}

\begin{lemma}[Five-cycles are commutators]
  \label{lem:circuits-five-cycles}
  \lean{Complexity.exists_fiveCycle_commutator,
    Complexity.every_fiveCycle_is_commutator}
  \leanok
  Some two $5$-cycles in $S_5$ have a $5$-cycle as their commutator, and
  every $5$-cycle in $S_5$ is the commutator $[a, b] = a b a^{-1} b^{-1}$ of
  two $5$-cycles $a$ and $b$.
\end{lemma}
\begin{proof}
  \leanok
  One explicit pair is checked by kernel computation, and the $5$-cycles form
  a single conjugacy class.
\end{proof}

\begin{theorem}[Barrington, finite form]
  \label{thm:circuits-barrington-finite}
  \lean{Complexity.Computes_formula_depth_four,
    Complexity.barrington_representation_depth_four,
    Complexity.barrington_quadratic_of_log_depth}
  \leanok
  \uses{def:circuits-formula, def:circuits-bp}
  For every $5$-cycle $\sigma \in S_5$, every Boolean formula of depth $d$ is
  computed through $\sigma$ by a width-$5$ permutation branching program of
  length at most $4^d$; in particular it is computed through some
  $\sigma \neq 1$ by such a program. A formula of depth at most
  $\lfloor \log_2 n \rfloor$, with $n \ge 1$, is therefore computed through
  some $\sigma \neq 1$ by a width-$5$ program of length at most $n^2$.
\end{theorem}
\begin{proof}
  \leanok
  \uses{lem:circuits-five-cycles}
  Induction on the formula, for all target $5$-cycles at once. The variable
  $x_i$ is the instruction $(i, 1, \sigma)$, the constant true is a single
  constant instruction, and false is the empty program. For a negation,
  compute the subformula through $\sigma^{-1}$ and fold a final
  multiplication by $\sigma$ into the last instruction, so the length does not
  grow. For a conjunction, write $\sigma = [a, b]$ with $5$-cycles $a$ and
  $b$ and concatenate programs for the two conjuncts through $a$ and $b$ and
  their inverses, four blocks of length at most $4^{d-1}$. A disjunction is a
  negated conjunction of negations, with the negations folded in the same
  way.
\end{proof}

\begin{theorem}[Constructive Barrington compiler]
  \label{thm:circuits-barrington-compiler}
  \lean{Complexity.barringtonCompile, Complexity.barringtonTargetBase,
    Complexity.barringtonTargetBase_spec, Complexity.barringtonCompile_computes,
    Complexity.barringtonCompile_length_le, Complexity.barringtonCompile_var_bound,
    Complexity.barringtonCompile_representation}
  \leanok
  \uses{def:circuits-formula, def:circuits-bp}
  An explicit computable compiler maps every formula $\varphi$ and target
  $\sigma \in S_5$ to a width-$5$ program of length at most
  $4^{\mathrm{depth}(\varphi)}$ which, whenever $\sigma$ is a $5$-cycle,
  computes $\varphi$ through $\sigma$. If every variable of $\varphi$ is at
  most $b$, then so is the variable read by every instruction. For the fixed
  $5$-cycle $\sigma_0 = [a_0, b_0]$, the commutator of two explicit
  $5$-cycles, the program evaluates to $\sigma_0 \neq 1$ when $\varphi$ is true
  and to the identity otherwise.
\end{theorem}
\begin{proof}
  \leanok
  \uses{lem:circuits-five-cycles}
  Follow the recursion of Theorem~\ref{thm:circuits-barrington-finite},
  writing each target $5$-cycle $\sigma$ as the commutator of the conjugates
  of $a_0$ and $b_0$ by the first permutation, in an explicit enumeration of
  $S_5$, that conjugates $\sigma_0$ to $\sigma$.
\end{proof}

\begin{lemma}[Programs to balanced formulas]
  \label{lem:circuits-bp-to-formula}
  \lean{Complexity.BP.decisionFormula, Complexity.BP.eval_decisionFormula_eq_true,
    Complexity.BP.depth_decisionFormula_le, Complexity.BP.vars_decisionFormula_lt}
  \leanok
  \uses{def:circuits-formula, def:circuits-bp}
  For a width-$w$ program $P$ of length $\ell$ and a point $x$, there is a
  formula that, under every assignment, is true exactly when the value of $P$
  moves $x$, and that has depth at most $(w + 1) \lceil \log_2 \ell \rceil + 2$,
  where $\lceil \log_2 \ell \rceil$ is read as $0$ for $\ell \le 1$. If every
  instruction of $P$ reads a variable below $n$, then every variable of the
  formula is below $n$.
\end{lemma}
\begin{proof}
  \leanok
  Split the program in half and compose the two halves through the $w$
  possible intermediate points.
\end{proof}

\begin{definition}[Formula and program families]
  \label{def:circuits-formula-nc1}
  \lean{Complexity.FixedArityFormulaFamily, Complexity.FixedArityBPFamily,
    Complexity.FormulaNC1, Complexity.Width5BP}
  \leanok
  \uses{def:circuits-formula, def:circuits-bp}
  $\cc{FormulaNC1}$ is the class of Boolean function families computed by a
  family of formulas $(\varphi_n)$ in which $\varphi_n$ reads only variables
  below $n$, is evaluated on $x \in \{0,1\}^n$ under the assignment sending
  $i < n$ to $x_i$ (and every other index to $0$), and has depth at most
  $c \lfloor \log_2 n \rfloor + c$ for a constant $c$. $\cc{Width5BP}$ is the
  class of Boolean function families decided by a family of width-$5$
  programs $P_n$ for $n \ge 1$, each reading only variables below $n$ and with
  a designated point $x_n$, together with an explicit answer at $n = 0$; on
  $x \in \{0,1\}^n$ the output is $1$ exactly when the value of $P_n$ under
  the same assignment moves $x_n$, and the length of $P_n$ is at most
  $C (n + 1)^k$ for constants $C, k$.
\end{definition}

\begin{theorem}[Barrington]
  \label{thm:circuits-barrington}
  \lean{Complexity.fixedArityFormulaFamily_exists_bp,
    Complexity.fixedArityBPFamily_exists_formula, Complexity.barrington_equivalence,
    Complexity.FormulaNC1OnTotalAssignments, Complexity.Width5BPOnTotalAssignments,
    Complexity.barrington_equivalence_onTotalAssignments}
  \leanok
  \uses{def:circuits-formula-nc1}
  $\cc{FormulaNC1} = \cc{Width5BP}$: every formula family witnessing
  membership in $\cc{FormulaNC1}$ has a program family witnessing membership in
  $\cc{Width5BP}$ that computes the same Boolean function family, and
  conversely. The same equality holds for the total-assignment variants, whose
  members assign to each $n$ a function $(\mathbb{N} \to \{0,1\}) \to \{0,1\}$:
  there the formula at $n$ may read any variable and has depth at most
  $c \lfloor \log_2 n \rfloor + c$, and the program at $n$ may read any
  variable, decides by whether it moves a designated point, and has length at
  most $C (n + 1)^k$ at every $n$, including $n = 0$.
\end{theorem}
\begin{proof}
  \leanok
  \uses{thm:circuits-barrington-compiler, lem:circuits-bp-to-formula}
  Compile $\varphi_n$ at the fixed $5$-cycle $\sigma_0$. The program evaluates
  to $\sigma_0$ or to the identity, and a $5$-cycle moves every point, so any
  designated point is moved exactly when $\varphi_n$ is true. Its length is at
  most $4^{c \lfloor \log_2 n \rfloor + c} \le 4^c n^{2c}$, and it reads only
  variables below $n$ (constant instructions read variable $0$). Conversely,
  a program of length at most $C (n + 1)^k$ has a balanced decision formula of
  depth at most $6 \lceil \log_2 (C (n + 1)^k) \rceil + 2$, which is
  logarithmic in $n$, with its variables below $n$.
\end{proof}

\begin{theorem}
  \label{thm:circuits-nc1-subset-width5bp}
  \lean{Complexity.NC1_subset_FormulaNC1, Complexity.NC1_subset_Width5BP}
  \leanok
  \uses{def:NC-AC, def:circuits-formula-nc1}
  $\cc{NC}^1 \subseteq \cc{FormulaNC1}$, and hence
  $\cc{NC}^1 \subseteq \cc{Width5BP}$.
\end{theorem}
\begin{proof}
  \leanok
  \uses{lem:circuits-output-formula, thm:circuits-barrington}
  Unfold the output of each circuit; the depth at most doubles, from
  $c \lfloor \log_2 n \rfloor + c$ to
  $2c \lfloor \log_2 n \rfloor + 2c$, and the variables stay below the arity.
  The polynomial size bound is not used.
\end{proof}

\begin{theorem}
  \label{thm:circuits-formula-nc1-subset-nc1}
  \uses{def:NC-AC, def:circuits-formula-nc1}
  $\cc{FormulaNC1} \subseteq \cc{NC}^1$. Consequently
  $\cc{NC}^1 = \cc{FormulaNC1} = \cc{Width5BP}$.
\end{theorem}
\begin{proof}
  \uses{thm:circuits-nc1-subset-width5bp, thm:circuits-barrington}
  A formula of depth $d$ has fewer than $2^{d+1}$ nodes, so logarithmic depth
  gives polynomial size. At each length $n \ge 1$, compile each node to a
  gate, pushing negations down by De Morgan's laws into input flags and
  building constants from an input and its negation; the stored answer at
  length $0$ is the value of $\varphi_0$.
\end{proof}

\begin{lemma}[Bit-level Barrington generator]
  \label{lem:circuits-barrington-codegen}
  \lean{Complexity.FormulaCode.encode, Complexity.FormulaCode.encode_injective,
    Complexity.BPCode.Program.encode, Complexity.BPCode.Program.encode_injective,
    Complexity.barringtonCompileCode, Complexity.barringtonCompileCode_encode,
    Complexity.barringtonCompileCode_spec,
    Complexity.length_barringtonCompileCode_encode_le}
  \leanok
  \uses{thm:circuits-barrington-compiler}
  There are injective codes for formulas (a unary token count followed by the
  postfix token stream) and for width-$5$ programs (a unary instruction count
  followed by the instructions, each a unary variable index and two seven-bit
  permutation ranks), and a total function $g$ on bit strings, sending every
  string that does not decode to a formula to the empty string, such that for
  every formula $\varphi$ of depth $d$, $g(\mathrm{code}(\varphi))$ is the
  code of, and decodes to, the program compiled from $\varphi$ at $\sigma_0$
  in Theorem~\ref{thm:circuits-barrington-compiler}, and
  $|g(\mathrm{code}(\varphi))| \le 4^d + 1 + 4^d (|\mathrm{code}(\varphi)| + 15)$.
\end{lemma}
\begin{proof}
  \leanok
  Run the compiler on the decoded formula and bound each serialized
  instruction by the code length of the formula.
\end{proof}

\begin{lemma}[The Barrington generator in small space]
  \label{lem:circuits-barrington-codegen-space}
  \uses{lem:circuits-barrington-codegen}
  A deterministic transducer computes $g(\mathrm{code}(\varphi))$ from
  $\mathrm{code}(\varphi)$ using
  $O(\mathrm{depth}(\varphi) + \log |\mathrm{code}(\varphi)|)$ work-tape
  cells.
\end{lemma}
\begin{proof}
  Emit the compiled program instruction by instruction: an instruction index
  is a path of length $\mathrm{depth}(\varphi)$ in the recursion tree, with a
  constant amount of permutation data per level, and the formula node at the
  current position is located by counting in the input code.
\end{proof}

\begin{theorem}[Uniform Barrington]
  \label{thm:circuits-uniform-barrington}
  \uses{def:circuits-formula-nc1, def:FL}
  \notready
  Call a formula family logspace-uniform if some function in $\cc{FL}$ maps
  $1^n$ to $\mathrm{code}(\varphi_n)$, and a program family logspace-uniform if
  some function in $\cc{FL}$ maps $1^n$ to the code of $(P_n, x_n)$. A Boolean
  function family has a logspace-uniform $\cc{FormulaNC1}$ witness if and only
  if it has a logspace-uniform $\cc{Width5BP}$ witness.
\end{theorem}
\begin{proof}
  \uses{thm:circuits-barrington, lem:circuits-barrington-codegen-space,
    lem:circuits-bp-to-formula}
  Forward: compose the formula generator with the small-space Barrington
  generator; the depth is $O(\log n)$, so the space is $O(\log n)$.
  Backward: generate the balanced decision formula of each program in
  logarithmic space.
\end{proof}

\section{Open directions}
\label{sec:circuits-open}

The following items are infrastructure rather than theorems.

\begin{itemize}
  \item \textbf{A circuit encoding designed for generation.} The current code
    format has a gate-count header, absolute wire references into subformulas
    of varying size, and no sharing. The log-space generator behind
    Lemma~\ref{lem:circuits-padded-code-fl} must therefore compute closed-form
    wire offsets for every fragment, and the uniform unrolling development is
    about 59,000 lines. A layered grid encoding, in which every wire reference
    is an affine function of the loop counters, would make future uniformity
    proofs (uniform $\cc{NC}$ and $\cc{AC}$, Theorem~\ref{thm:circuits-uniform-barrington},
    finer notions such as DLOGTIME uniformity) much shorter. It should come
    with an $\cc{FL}$ translation into the existing codec, so that
    $\cc{UniformPPoly}$ does not change.
  \item \textbf{One representation of circuits and formulas.} The library has
    about fourteen parallel representations: typed circuits, raw gate lists,
    general, monotone and negation-normal unbounded formulas, two CNF/DNF
    types, two decision-tree types, branching programs, and several family
    wrappers. A single formula type with a substitution operation (a monadic
    bind), and circuits built so that acyclicity holds by typing, would let
    restriction, unfolding and normalization lemmas be proved once.
\end{itemize}
